%% transmon-benchmark.23 -- v22 revised for reading flow, 2026-10-23: paragraphs,
%% a float and one section moved; terms pinned at first use. No number, result,
%% claim scope, title, author or DOI changed; no solver was run. See
%% changelog.md, provenance.md and validation.md. Earlier versions are unchanged.

\documentclass{anvil-paper}
\usepackage{xurl}
\usepackage{array}
\newcommand{\Cisl}{C_{\mathrm{isl}}}

%% --- URL line-breaking: allow breaks at hyphens and add stretch so the
%% bibliography stays inside the measure (render-gate overfull check).
\expandafter\def\expandafter\UrlBreaks\expandafter{\UrlBreaks\do\-}
\Urlmuskip=0mu plus 1mu
%% Float placement: let text share pages with floats.
\renewcommand{\topfraction}{0.92}
\renewcommand{\bottomfraction}{0.6}
\renewcommand{\textfraction}{0.07}
\renewcommand{\floatpagefraction}{0.85}
\setcounter{topnumber}{3}
\setcounter{totalnumber}{4}
%% Compact bibliography.
\renewcommand{\bibfont}{\fontsize{8}{8.8}\selectfont}
\setlength{\bibsep}{0pt plus 1pt}
\setlength{\emergencystretch}{1.5em}

%% --- Draft-phase helper macros -------------------------------------------
%% Figure placeholder: uses the rendered figure if pub-figures has produced
%% it, otherwise a framed placeholder naming the source script.
\newcommand{\anvilfig}[3][0.98]{%
  \IfFileExists{figures/#2}%
    {\includegraphics[width=#1\linewidth]{figures/#2}}%
    {\fbox{\parbox{#1\linewidth}{\centering\vspace{5em}%
      \texttt{figures/#2}\\[0.5em]%
      {\small placeholder: render via \texttt{figures/src/#3}}%
      \vspace{5em}}}}}

\title{An Independent Finite-Element Cross-Check\\of the DeviceLayout.jl Transmon}
\author{Robb Walters\thanks{Corresponding author: \href{mailto:robb@2amlogic.com}{robb@2amlogic.com}}\\
  {\normalsize 2AM Logic}
  \and Crutcher Dunnavant\\
  {\normalsize ZSpaceLabs.ai}}
\date{October 2026\\[0.5ex]
  {\normalsize DOI: \href{https://doi.org/10.5281/zenodo.23580545}{10.5281/zenodo.23580545}}}

\begin{document}
\maketitle

\begin{abstract}
DeviceLayout.jl's SingleTransmon example publishes qubit and resonator
frequencies for an initial design and a final design optimized for $3.0$
and $4.0$\,GHz, computed with Palace on one mesh. We find them
discretization-limited. Palace reproduces them within $0.7\%$ at the
example's own discretization, and the independent solver geode-fem agrees
on shared meshes to $2\times10^{-4}\%$ (qubit) and $0.03\%$
(resonator). All frequencies are linear eigenmodes. Palace's
still-rising adaptive refinement puts the initial and final designs' qubit
frequencies $7.0\%$ and $8.0\%$ above their published values without
extrapolation. A two-sided bound on the initial design's island
capacitance, read through a circuit model, caps the qubit at
$4.477$\,GHz ($+8.1\%$). Uniform refinement converges slowly because of
the singular metal edges. Extrapolated assuming error $\propto h$, the
fastest rate those edges allow on uniform meshes, two refinement families put both
designs' qubit and resonator frequencies $7.3$--$9.6\%$ and
$2.7$--$3.3\%$ above their published values, estimates that lean low and
rest on Palace's finer levels. At these frequencies the final design needs
a capacitor $19$--$20\%$ longer than published. Newton's method in
DeviceLayout's variables meets the example's criterion in 3 or 4 Palace
runs with any of three Jacobians, against 10 to 12 for PRIMA; geode-fem's
shape derivative gives no advantage here.
\end{abstract}

\section{Introduction}
\label{sec:intro}

A transmon qubit is an electrical oscillator: a metal island, whose
capacitance to its surroundings is the oscillator's capacitor, connected
to ground through a Josephson junction, which acts as its inductor. Its
frequency, which we call the \emph{qubit frequency}, is the linear
eigenfrequency of the junction's (plasma) mode.\footnote{Not the
transmon's 0--1 transition $f_{01}$, which lies below the plasma frequency
by about the charging energy $E_C/h$ (Section~\ref{sec:junction}).} It is
set mostly by that capacitance, which a three-dimensional field solve
computes, and the readout resonator coupled to it is sized the same way. Designers lengthen or shorten the island and the resonator until
the computed frequencies sit on their targets, so if a solve reads a
frequency $8\%$ low, the design it steers to has, in the same device
model, a frequency about $8\%$ above its target.

DeviceLayout.jl \citep{devicelayout} builds such a layout and its mesh, and
Palace \citep{palace}, an MFEM-based finite-element solver \citep{mfem},
computes the eigenmodes; tools such as QDesignOptimizer
\citep{eriksson2025automated} steer the same loop with analytic models.
DeviceLayout's SingleTransmon example
\citep{devicelayout-blog,devicelayout-example} runs the loop end to end. It
publishes an initial transmon-plus-resonator design at $4.14$ and
$5.591$\,GHz and a design optimized for targets of $3.0$ and $4.0$\,GHz, at
$3.005$ and $3.982$\,GHz, which we call the \emph{final} design; a
derivative-free optimizer found it in nine Palace runs (one Palace
eigensolve each, the unit in which this paper counts optimization cost,
Section~\ref{sec:reopt}). Like every frequency here, these are linear
eigenfrequencies.

Two solvers appear. Palace is the one the example uses; the second is
GEODE-FEM (``geode-fem''), a Rust finite-element solver on the Burn tensor
framework \citep{burn}. It solves the same
$H(\mathrm{curl})$ eigenproblem as Palace with a separate implementation,
and it computes design derivatives, among them the shape derivative of an
eigenfrequency. Given the same input mesh, it reproduces Palace's modes
(Section~\ref{sec:agreement}), and on the $h/2$ meshes of
Section~\ref{sec:crosscheck} the two codes solve the same discrete problem
apart from the junction term. This gives an independent
same-mesh cross-check of the Palace/DeviceLayout.jl transmon stack,
complementing the earlier checks of Palace by SQDMetal
\citep{sommers2025open} and \citet{ye2025electromagnetic}.
The published frequencies come from one discretization, and the example's
authors say what that implies. They run Palace at second order ``to reduce
sensitivity to the discretization'' and note that this ``does not mean that
the eigenfrequencies found by Palace will have converged to their true
values'' \citep{devicelayout-example}. This paper measures
how far from converged the published frequencies are. Our
main result is that the published frequencies are discretization-limited.
Palace's adaptive refinement puts the qubit frequencies of the initial and
final designs $7.0\%$ and $8.0\%$ above their published values without
extrapolation, and its
uniform ladders, in two refinement families (two ways of subdividing the
same starting mesh), extrapolate to $7.3$--$9.6\%$ above them;
every uniform level we could afford lies below those estimates. geode-fem
reproduces every level of those ladders it can reach, the two coarsest
meshes with linear elements, to $10^{-4}\%$ on the qubit, and a
two-sided bound on the island capacitance, read through a circuit model
(the \emph{static route}, Section~\ref{sec:static}), is consistent with
that shift and caps the initial qubit under that model. The extrapolation
rule was chosen after that capacitance route had been computed, so their
agreement is a consistency check, not a blind prediction.
Section~\ref{sec:discussion} explains why the same-mesh agreement places
the shift in the mesh, not in either solver or its arithmetic.

Five words sort the evidence, and Table~\ref{tab:tiers}, in the
Discussion, sorts every result by the tiers they name. A \emph{bound}
holds by theorem and is given only for the island capacitance (the
\emph{Bounded} tier); a \emph{cap} is a frequency limit that follows from
a bound only under the circuit model (\emph{Model-dependent}); an
\emph{estimate} is an extrapolation, quoted with its \emph{band}, the
range of the extrapolations, which is neither a bound nor a cap
(\emph{Estimated}); and a \emph{computed} value is a solve on a stated
discretization, with no extrapolation (\emph{Computed}).

Our contributions are:
\begin{itemize}
  \item \textbf{A refinement study of a published design example}, which
    puts its qubit frequencies $7.3$--$9.6\%$ above their published values, with a
    second refinement family and adaptive runs on both designs checking
    the qubit estimates; the finer levels and the extrapolations are
    Palace's (Section~\ref{sec:crosscheck}).
  \item \textbf{Two-sided capacitance bounds with the sapphire tensor.}
    An equilibrated-flux lower bound and conforming upper bounds bracket
    the island capacitance $\Cisl$ of the unmodified device, without the
    junction capacitance, in $[79.56, 79.77]$\,fF with the box walls held
    at ground, as in the eigenproblem; under a two-node circuit model the
    lower bound caps the initial qubit at $4.477$\,GHz
    (Section~\ref{sec:static}).
  \item \textbf{A re-optimized design, with controls.} At the extrapolated
    frequencies a design that reaches the final design's targets, $3.0$
    and $4.0$\,GHz, needs $c \approx 1451$--$1466\,\mu$m against the
    published $1217.678\,\mu$m. Against
    the example's own objective, Newton's method in DeviceLayout's
    variables needs 3 or 4 Palace runs with any of three Jacobians and
    PRIMA 10 to 12. geode-fem's shape derivative is exact but gives no
    advantage on this two-parameter problem; its one-time calibration of
    the resonator model replaces one Palace probe run
    (Sections~\ref{sec:reopt} and~\ref{sec:discussion}).
  \item \textbf{Same-input-mesh agreement with Palace}, with the
    $17.49$\,GHz mode that both codes share traced to the port model
    (Section~\ref{sec:agreement}).
\end{itemize}

\section{Related Work}
\label{sec:related}

\paragraph{Cross-checks of Palace.}
SQDMetal \citep{sommers2025open} cross-checks a Palace workflow against
COMSOL and Ansys with mesh-convergence studies, and
\citet{ye2025electromagnetic} compare a Palace workflow with cryogenic
measurement, matching resonator frequencies to within $0.3\%$ with a mesh
refined at the CPW gaps and fourth-order elements, a setting their
convergence study puts within $\pm0.3\%$ of extrapolated values. Both cross-checks use models built
independently for each tool. Our comparison instead holds the input mesh
and the junction model fixed and changes the solver implementation, and it
adds a refinement study of a published design: code-to-code comparison and
solution verification \citep{oberkampf2010verification}, not validation
against measurement. As of October 2026, neither Palace nor DeviceLayout.jl has a
methods publication; we cite their repositories and documentation.

\paragraph{Design sensitivities and qubit-design optimization.}
Adjoint sensitivity analysis of frequency-domain electromagnetic solvers is
an established field \citep{nikolova2004adjoint}, including shape adjoints
for Maxwell cavity eigenvalues \citep{akcelik2005adjoint}; Palace exposes no
derivatives.
In qubit design, SQcircuit \citep{rajabzadeh2023analysis} and
SuperGrad \citep{wang2025supergrad} differentiate circuit Hamiltonians,
SQuADDS \citep{shanto2024squadds} looks designs up in what its authors
call a ``validated design database'', and \citet{eun2023shape} shape-optimize transmon pads with FEM
and a derivative-free optimizer.
DeviceLayout's example optimizes with PRIMA \citep{zhang2023prima}, a
reference implementation of Powell's model-based derivative-free methods;
with bounds only it runs BOBYQA \citep{powell2009bobyqa}.

\paragraph{Two-sided capacitance bounds and circuit models.}
Bracketing a capacitance between complementary energy principles is
classical \citep{specogna2011complementary}. Our lower bound uses the
hypercircle idea of \citet{prager1947approximations}, which turns any
field that nearly balances the charges into a lower bound once its
imbalance is paid for, with Raviart--Thomas
fluxes \citep{raviart1977mixed} equilibrated patch by patch
\citep{braess2007equilibrated,ern2015polynomial}. The method is not new; we
apply it to an island capacitance with a floating conductor on a
production qubit layout. A web search in October 2026 of the
complementary-bound and qubit-design literature (queries such as
``guaranteed capacitance bounds qubit'') found no two-sided capacitance
bound in the qubit-design literature. The two-node circuit of Section~\ref{sec:junction} is a minimal instance of
the quasi-lumped route \citep{minev2021circuit}, which partitions a device
into lumped or quasi-distributed cells, simulates each and reduces the
composite to coupled subsystems; the energy-participation route
\citep{minev2021energy} instead quantizes the distributed eigenmodes.

\section{Device, Models and Solvers}
\label{sec:setup}

\subsection{Device, meshes and materials}
\label{sec:geometry}

The device is DeviceLayout.jl v1.15.0's SingleTransmon example
\citep{devicelayout}: a transmon \citep{koch2007charge} with a meander readout resonator and a
feedline on sapphire, in a vacuum box (Figure~\ref{fig:geometry}). Its two
design parameters are the transmon's \texttt{capacitor\_length} $c$, the
length of the island, and the resonator's \texttt{total\_length} $L$, the
length of the meander (Figure~\ref{fig:geometry}a); the published initial design has
$c = 620\,\mu$m and $L = 5000\,\mu$m. The sapphire permittivity is
$\mathrm{diag}(9.3, 9.3, 11.5)$ with the optic axis normal to the chip;
the configuration rotates the material axes about that normal, which
leaves a uniaxial tensor unchanged, so the lab tensor is diagonal. Where a
scalar model is used it is the trace average $\bar\varepsilon = 10.03$.
PEC (perfect-conductor) conditions hold on the metal sheets and on the box
walls, which therefore sit at ground potential; the readout ports are left
open in every eigenproblem except the DeviceLayout-configuration runs of
Section~\ref{sec:crosscheck}.

\begin{figure}[htbp]
  \centering
  \anvilfig{fig1-geometry.pdf}{fig1\_geometry.py}
  \caption{Benchmark geometry, drawn from the SHA-256-pinned MSH~4.1
  input mesh of the initial design (the fixture; 22{,}684 nodes,
  133{,}314 tetrahedra). (a) Top view of the three metal conductors of
  Section~\ref{sec:static}: the feedline (teal), the qubit island (blue,
  as in (b) and (d)), and the ground plane (gray) with the meander
  resonator (its outline orange), whose length is $L$; the junction
  (purple, as in (b) and as $L_J$ in (d)) joins the island to ground, squares mark the two
  readout ports, and the dashed outline is the box wall. Each label has the
  colour of the object it names, except the $k$, Vacuum, Sapphire and
  Metal-sheets labels; the design parameters are annotated in green. $c$ is the island length. When $c$
  changes, the shaded strip stretches and everything above it moves up by
  $\Delta c$ (Section~\ref{sec:reopt}; bands from \texttt{fd\_initial.toml}). (b) The $1\times1\,\mu$m junction
  \texttt{lumped\_element} sheet with its four-triangle triangulation.
  (c) Material cross-section with the box walls (dashed), schematic and not
  to scale. (d) The two-node circuit of Section~\ref{sec:junction}: the
  island's capacitance $C_\Sigma$ across the junction inductance $L_J$ and
  the series inductance $L_s$ of its current path, coupled with
  coefficient $k$ to the resonator.}
  \label{fig:geometry}
\end{figure}

Table~\ref{tab:meshes} lists the meshes. Element order $p$ is the
polynomial degree of the finite elements ($p = 1$: linear, $p = 2$:
quadratic); \emph{straight-sided} and \emph{curved} describe the geometry
of the mesh, whose element faces are flat or follow curved surfaces. The
\emph{fixture} is a straight-sided mesh of the initial geometry written by
an earlier DeviceLayout.jl v1.15.0 driver,
\path{reference/julia/generate_transmon_fixture.jl}, pinned by SHA-256;
its levels are $h_F$, $h_F/2$ and $h_F/4$. Refining it leaves its
polygonal surfaces unchanged, which is what the capacitance bounds of
Section~\ref{sec:static} need. For the cross-check and the
re-optimization, DeviceLayout meshes each geometry with the two parameters
exposed, straight-sided at level $h$ for the ladders and curved for the
runs that reproduce its discretization; curving the faces changes the
qubit by at most $0.0001\%$ and the resonator by about $0.05\%$. On the
initial geometry $h_F$ and $h$ give the same Palace qubit ($3.4527$
against $3.4520$\,GHz, $0.02\%$). A refinement \emph{family} is one way of
subdividing $h$ uniformly: family~a with gmsh, family~b with Palace; a
\emph{ladder} is one geometry's series of uniformly refined levels, from
$h$ down, at each element order.

\begin{table}[htbp]
  \centering
  \caption{The meshes. ``Initial'' and ``final'' are the two published
  designs; ``step 2'' is the re-optimized design of
  Section~\ref{sec:reopt}.}
  \label{tab:meshes}
  \footnotesize
  \begin{tabular}{@{}>{\raggedright\arraybackslash}p{0.16\linewidth}>{\raggedright\arraybackslash}p{0.12\linewidth}>{\raggedright\arraybackslash}p{0.27\linewidth}>{\raggedright\arraybackslash}p{0.11\linewidth}>{\raggedright\arraybackslash}p{0.22\linewidth}@{}}
    \toprule
    Mesh & Designs & How made & Levels & Used for \\
    \midrule
    Fixture & initial & earlier DeviceLayout driver, straight-sided; refined uniformly, or adaptively by geode-fem & $h_F$, $h_F/2$, $h_F/4$; adapted & same-mesh agreement (\S\ref{sec:agreement}); capacitance bounds and cap (\S\ref{sec:static}) \\
    DeviceLayout $h$ & every design & DeviceLayout, straight-sided & $h$ & ladders (\S\ref{sec:crosscheck}); geode-fem solves and derivatives (\S\ref{sec:reopt}) \\
    Family a & every design & gmsh splits each tetrahedron of $h$ into 8 & $h/2$, $h/4$ & ladders and estimates \\
    Family b & initial, final, step 2 & Palace's uniform refinement of $h$ & $h/2$, $h/4$ & a second estimate \\
    Adaptive & initial, final & Palace adaptive refinement from $h$ at $p = 2$ & passes 1--5 & check of the qubit estimates \\
    DeviceLayout's discretization & every design & DeviceLayout, curved; Palace at $p = 2$ & one & reproduction; the Newton objective $F$ (\S\ref{sec:reopt}) \\
    \bottomrule
  \end{tabular}
\end{table}

\subsection{Junction and circuit models}
\label{sec:junction}

The Josephson junction is a linear lumped element, $L_J = 14.860$\,nH in
parallel with $C_J = 5.5$\,fF, on the four-triangle sheet $\Gamma$ of
Figure~\ref{fig:geometry}b, as in DeviceLayout's Palace configuration. With
port length $\ell$ and width $w$ the uniform lumped-port construction adds
\begin{equation}
  K_{\mathrm{port}} = \frac{\ell}{w\,L_J}\, S_\Gamma
  \qquad\text{and}\qquad
  M_{\mathrm{port}} = \frac{C_J\,\ell}{w}\, S_\Gamma
  \label{eq:port}
\end{equation}
to the pencil $(K + K_{\mathrm{port}})x = \omega^2(M + M_{\mathrm{port}})x$,
where $S_\Gamma$ is the isotropic surface matrix on $\Gamma$. Palace's
\texttt{LumpedPort} adds the same terms with the same scalings. geode-fem's
rank-1 port-voltage shunt replaces $S_\Gamma$ by $b\,b^{\mathsf T}$, with
$b_i = (1/w)\int_\Gamma N_i\cdot\hat e\,dS$ the port-voltage functional; on
the fixture it moves the qubit by $6\times10^{-5}\,\%$, keeps the other
physical modes to four decimals, removes the
artifact mode of Section~\ref{sec:agreement}, and is the junction model of
every geode-fem eigensolve in Sections~\ref{sec:crosscheck}
and~\ref{sec:reopt}.

The electrostatic model splits the metal into three conductors: ground
with the resonator, the floating feedline, and the island. Its island
capacitance with the feedline floating and neutral,
$\Cisl = C_{ii} - C_{if}^2/C_{ff}$, excludes the junction. In the
convention of \citet{koch2007charge} the transmon's total capacitance is
$C_\Sigma = \Cisl + C_J$, and it sets the charging energy $E_C =
e^2/2C_\Sigma$. A two-node circuit (Figure~\ref{fig:geometry}d) connects
it to the eigenproblem: a bare
qubit $1/(2\pi\sqrt{(L_J + L_s)C_\Sigma})$, with $L_s$ the series
inductance of the junction's current path, coupled with coefficient $k$
to the resonator. \emph{Bare} is the qubit's own LC frequency;
\emph{dressed} is its frequency after the coupling to the resonator
pushes it, which is what the eigenproblem computes. The same circuit is
used both ways: forward from a
bound on $C_\Sigma$ to a dressed qubit, and backward from a dressed pair to
$C_\Sigma$. At $L_s = k = 0$ this is the plasma frequency
$\sqrt{8E_JE_C}/h$, and the transmon's $f_{01} \approx \sqrt{8E_JE_C}/h -
E_C/h$ lies below it. On the fixture $L_s$ and $k$ were fitted on three and
two discretizations (\path{benchmarks/transmon_loop_inductance/},
\path{benchmarks/transmon_coupling/}); extrapolated, they give fitted ranges of $0.113$--$0.125$\,nH
and $0.0227$--$0.0310$, which are estimates.

\subsection{geode-fem}
\label{sec:architecture}

geode-fem assembles first-order N\'ed\'elec \citep{nedelec1980} and scalar
P1/P2 element operators as batched $[n_{\mathrm{elem}},k,k]$ tensor
operations on Burn. Its eigensolver is shift-invert Lanczos \citep{arpack}
on a real symmetric pencil, factorized with faer \citep{kazdadi2026faer}. The two codes share the
formulation, the lumped junction port, the materials and the input mesh,
and geode-fem follows a specification of Palace's operators
(Section~\ref{sec:built}); assembly, the eigensolver and the linear algebra
are separate implementations.

The eigenfrequency derivative used in Section~\ref{sec:reopt} needs no
adjoint solve, because the pencil is symmetric: for a simple eigenpair
$(\lambda, x)$ with $x^{\mathsf T}Mx = 1$,
\begin{equation}
  \frac{d\lambda}{d\theta} \;=\;
  x^{\mathsf T}\Bigl(\frac{\partial K}{\partial\theta}
  - \lambda\,\frac{\partial M}{\partial\theta}\Bigr)x ,
  \label{eq:hf-eig}
\end{equation}
the classical eigenvalue derivative of a symmetric generalized eigenproblem
\citep{foxkapoor1968rates}. Here $K$ and $M$ include the junction terms of
Eq.~\eqref{eq:port}; the junction does not move under the capacitor-length
motion of Section~\ref{sec:reopt}, so those terms do not contribute.
geode-fem evaluates $\partial K/\partial\theta$ and $\partial
M/\partial\theta$ along a node velocity $\partial X/\partial\theta$ with a
hand-written forward-mode dual-number twin of its N\'ed\'elec kernel
(tested against it to $10^{-10}$); Burn's autodiff tape is not used.

\section{Same-Input-Mesh Agreement With Palace}
\label{sec:agreement}

Both codes read the fixture with the same materials and junction model at
$p = 1$ (Table~\ref{tab:agreement}). Palace's configuration schema
defines \texttt{Target} as the frequency above which its eigensolver
searches \citep{palace}, so the six modes above $4.5$\,GHz come from a run
with \texttt{Target}~$=4.5$\,GHz and the qubit from one with
\texttt{Target}~$=2$\,GHz. All six agree to $\leq 0.033\%$, and the qubit mode to $5\times10^{-5}\,\%$, with port
participation $0.9942$ in geode-fem and port EPR $0.99424$ in Palace, the
energy-participation ratio of the junction, the fraction of the mode's
inductive energy stored in it \citep{minev2021energy}. Both codes use the
isotropic sheet term of Eq.~\eqref{eq:port} here.

The two codes read the same input mesh but do not solve the same unknowns.
Palace cracks each metal sheet into two sides by duplicating its vertices,
which leaves the PEC problem unchanged, and where a sheet is one element
wide, almost all on the resonator and feedline sheets, it also bisects the
elements on both sides (``local bisection for under-resolved interior
boundaries'' in its log; \texttt{palace/utils/geodata.cpp}). Only the
bisection changes the unknowns (Table~\ref{tab:agreement}).

\begin{table}[htbp]
  \centering
  \caption{Eigenmode agreement from the identical input mesh (22{,}684
  nodes, 133{,}314 tetrahedra), $p = 1$ in both solvers, junction as
  the lumped $L_J \parallel C_J$ shunt. Counted the same way, as all
  first-order N\'ed\'elec unknowns (edges) including those on PEC sheets,
  geode-fem's mesh has $156{,}863$ ($133{,}108$ interior) and Palace's has
  $183{,}891$, after it cracks the metal sheets and adds $4{,}913$ elements
  on them by bisection, triggered by $1{,}007$ edges, $993$ of them on
  the resonator's ground plane and the feedline. Both
  codes are printed to geode-fem's archived three decimals (six for the
  qubit); $\Delta$ is the archived relative difference, computed from
  Palace's full-precision value. Sources:
  \texttt{benchmarks/transmon\_eigen/results.toml} and
  \texttt{benchmarks/transmon\_eigen/palace\_crack/}.}
  \label{tab:agreement}
  \footnotesize
  \begin{tabular}{@{}llccc@{}}
    \toprule
    Mode & & geode-fem (GHz) & Palace (GHz) & $\Delta$ (\%) \\
    \midrule
    qubit & transmon qubit & 3.452751 & 3.452749 & 0.00005 \\
    resonator & readout resonator & 5.153 & 5.151 & 0.032 \\
    mode 2 & box/cavity & 15.465 & 15.461 & 0.029 \\
    port sheet & port-sheet artifact & 17.490 & 17.490 & 0.001 \\
    mode 4 & box/cavity & 18.693 & 18.692 & 0.007 \\
    mode 5 & box/cavity & 20.703 & 20.698 & 0.026 \\
    mode 6 & box/cavity & 26.088 & 26.081 & 0.027 \\
    \bottomrule
  \end{tabular}
\end{table}

The $17.49$\,GHz mode is a spurious mode of the junction-sheet model,
shared by both codes: Palace gives it a port EPR of $2.49\times10^{-8}$,
and geode-fem's rank-1 shunt removes it.\footnote{The junction patch has
one free node whose hat-function gradient carries no port voltage but
stores energy in $S_\Gamma$, with Rayleigh quotient $1/(L_JC_J)$, or
$f_{LC} = 17.60$\,GHz, lowered by the adjacent volume mass. This
mechanism is our inference, supported by a synthetic test.} The fixture
is coarse: Palace's own refinements of it move the qubit from
$3.4527$ to $3.9528$\,GHz at $h_F/2$ ($+14.5\%$) and to $4.1721$\,GHz at
$p = 2$ on $h_F$ ($+20.8\%$), and the resonator by $+6.1\%$ at $h_F/2$, so
the agreement above is a code-to-code comparison on one discretization.

\section{The DeviceLayout Cross-Check}
\label{sec:crosscheck}

\paragraph{The published designs.}
The DeviceLayout.jl blog \citep{devicelayout-blog} reports the initial
design ($c = 620\,\mu$m, $L = 5000\,\mu$m) at $[4.14, 5.591]$\,GHz and
the final design ($c = 1217.678\,\mu$m, $L = 7243.955\,\mu$m) at
$[3.005, 3.982]$\,GHz, reached by PRIMA in nine Palace runs ``within 1\% of
our target frequencies''. The example source \citep{devicelayout-example}
sets the targets to $[3.0, 4.0]$\,GHz, works in the variables $x_1 =
\sqrt{620\,\mu\mathrm{m}/c}$ and $x_2 = 5000\,\mu\mathrm{m}/L$, and stops
when $\mathrm{err} = \frac12\sum_m (f_m - t_m)^2/t_m^2 < 10^{-4}$. Its
frequencies are computed with Palace on DeviceLayout's curved mesh at
$p = 2$, with a first-order absorbing far field and two
$50\,\Omega$ readout ports; we call this \emph{DeviceLayout's
discretization}.

\paragraph{Reproduction at DeviceLayout's discretization.}
With the example's configuration file, Palace at DeviceLayout's
discretization gives $[4.1683, 5.6132]$\,GHz for the initial design
and $[3.0164, 3.9956]$\,GHz for the final one, $0.68\%$ and $0.40\%$
above the published initial pair and $0.38\%$ and $0.34\%$ above the
final pair. The same cells with the PEC-box model used everywhere else
differ from the configuration-file runs by at most $0.003\%$. The
published runs' Palace version, mesher build and platform are not
recorded (ours are in Section~\ref{sec:repro}); every published value lies
below ours, as a slightly coarser published mesh would make it, but which
difference causes the $0.34$--$0.68\%$ gap is not diagnosed.

\paragraph{The ladder and its rate.}
For each geometry, Palace runs at $p = 1$ on DeviceLayout's straight-sided
mesh $h$ and its uniform refinements $h/2$ and $h/4$ (family~a, up to
$12.6$ million unknowns), and at $p = 2$ on $h$ and $h/2$;
geode-fem runs at $p = 1$ on $h$ and $h/2$, on the same meshes. Every
computed series rises under refinement (Figure~\ref{fig:crosscheck}); this
is an observation, not a guarantee. We write the discretization error as
$\propto h^{r}$ and call $r$ the \emph{convergence rate}. At the edge of a
zero-thickness metal sheet the field behaves as $E \sim r_e^{-1/2}$ in the
distance $r_e$ from the edge \citep{meixner1972behavior},
so near the edge the eigenfunctions lie in $H^{s}$ only for $s < 1/2$, and
on uniform meshes the eigenvalue error decays as $h^{2s}$ at any element
order \citep{boffi2010finite}: $r < 1$. The $p = 1$
ladders show pre-asymptotic rates of $0.61$--$0.62$ (qubit) and $0.70$--$0.73$
(resonator).

\paragraph{Extrapolation rule.}
The \emph{estimate} is the mean of two extrapolations at rate $r = 1$ in the
frequency, $2f_2 - f_1$: one from the $p = 1$ pair ($h/2$, $h/4$) and one
from the $p = 2$ pair ($h$, $h/2$) (Figure~\ref{fig:crosscheck}, small open
markers).\footnote{The rule is Richardson extrapolation
\citep{richardson1911approximate} at the asymptotic rate rather than the
observed one; Roache's grid convergence index \citep{roache1994perspective}
with a factor of $1.25$, on the initial qubit's $p = 1$ ladder, is centred
on $4.155$\,GHz with a half-width of $0.65$\,GHz, which reaches
$4.81$\,GHz, above the circuit-model cap.} $r = 1$ is the fastest rate the theory
allows on uniform meshes, so the rule makes the smallest correction
consistent with it, and its estimates lean low; for the $p = 2$ pair the
rate is taken from theory, because two levels give no observed rate. The
\emph{band} is the minimum and maximum of six extrapolations, the same two
pairs taken three ways (values for the initial qubit):
\begin{itemize}\setlength{\itemsep}{0pt}
  \item at $r = 1$ in the frequency, the two values of the estimate
    ($4.432$ and $4.455$\,GHz);
  \item at $r = 1$ in the eigenvalue variable $\lambda\propto f^2$,
    $\sqrt{2f_2^2 - f_1^2}$, which lie lower for a series that rises
    under refinement ($4.415$ and $4.451$\,GHz); the $p = 1$ value sets
    the low end, except on family-b resonators, where the $p = 2$ value
    does;
  \item at the observed rate of the $p = 1$ ladder's three levels
    ($0.62$ and $0.61$ for the family-a qubits), borrowed by the $p = 2$
    pair on the assumption that both element orders share the
    pre-asymptotic rate, which lie higher ($4.677$ and $4.582$\,GHz); the
    $p = 1$ value sets the high end.
\end{itemize}
The band is an extrapolation range, not a bound. For the initial qubit,
the circuit-model cap of Section~\ref{sec:static} trims it at
$4.477$\,GHz. The rule was fixed after that static route was known
(Section~\ref{sec:intro}); extrapolation at the observed rate alone would
put the initial qubit at $4.63$\,GHz, above the cap, so the observed-rate
values enter only as the band's high side.

\begin{figure}[htbp]
  \centering
  \anvilfig{fig2-crosscheck.pdf}{fig2\_crosscheck.py}
  \caption{Resolution ladders of the two published DeviceLayout.jl
  SingleTransmon designs: Palace at $p = 1$ and $p = 2$ on DeviceLayout's mesh
  $h$ and its family-a refinements, and geode-fem at $p = 1$ on the same
  meshes (crosses); the open diamond, just right of $h$, is Palace at
  DeviceLayout's own discretization, that of the published values
  (dashed). The ``$h\to0$'' slots are categorical: the estimate with its
  band for the gmsh refinement family (a, labelled) and for Palace's
  refinement (b). The small open markers in slot a are the $r = 1$
  extrapolations of the $p = 1$ pair ($h/2$, $h/4$) and the $p = 2$ pair
  ($h$, $h/2$), joined to the finer level of each pair; the estimate is
  their mean. In (a), the static route of Section~\ref{sec:static}
  (natural box; the grounded bound lowers its low end by about $0.3$\,MHz)
  and the circuit-model cap, which trims the untrimmed
  bands drawn. Sources:
  \texttt{benchmarks/transmon\_devicelayout\_crosscheck/results.toml} and
  \texttt{benchmarks/transmon\_devicelayout\_controls/results.toml}.}
  \label{fig:crosscheck}
\end{figure}

\begin{table}[htbp]
  \centering
  \caption{Published frequencies, Palace at DeviceLayout's discretization
  (DL), the finest pass of Palace's adaptive run (a computed value, still
  rising), and the extrapolated estimates with their bands, in GHz, for
  the two refinement families of DeviceLayout's mesh $h$ (a: gmsh's
  eight-way split; b: Palace's uniform refinement). Percentages are
  relative to the published value. Sources:
  \texttt{benchmarks/transmon\_devicelayout\_crosscheck/results.toml}
  (family a) and \texttt{benchmarks/transmon\_devicelayout\_controls/results.toml}
  (adaptive, family b).}
  \label{tab:crosscheck}
  \footnotesize
  \setlength{\tabcolsep}{3pt}
  \begin{tabular}{@{}lccc@{\hspace{6pt}}ll@{}}
    \toprule
    & & & Adaptive, & \multicolumn{2}{c}{Extrapolated estimate [band]} \\
    \cmidrule(l){5-6}
    & Published & DL & finest pass & Family a & Family b \\
    \midrule
    Initial qubit & 4.14 & 4.1683 & 4.4297 ($+7.0\%$) & 4.444 ($+7.3\%$) $[4.415, 4.677]^{\ast}$ & 4.472 ($+8.0\%$) $[4.458, 4.510]^{\ast}$ \\
    Initial resonator & 5.591 & 5.6132 & 5.6986 ($+1.9\%$) & 5.740 ($+2.7\%$) $[5.723, 5.817]$ & 5.776 ($+3.3\%$) $[5.768, 5.800]$ \\
    Final qubit & 3.005 & 3.0164 & 3.2468 ($+8.0\%$) & 3.265 ($+8.7\%$) $[3.236, 3.471]$ & 3.292 ($+9.6\%$) $[3.278, 3.331]$ \\
    Final resonator & 3.982 & 3.9956 & 4.0552 ($+1.8\%$) & 4.088 ($+2.7\%$) $[4.077, 4.139]$ & 4.113 ($+3.3\%$) $[4.106, 4.125]$ \\
    \bottomrule
  \end{tabular}

  \smallskip
  {\footnotesize $^{\ast}$Under the circuit model of
  Section~\ref{sec:static}, the cap of $4.477$\,GHz trims the bands to
  $[4.415, 4.477]$ and $[4.458, 4.477]$.}
\end{table}

\paragraph{Result.}
The qubit frequency extrapolates to $4.444$\,GHz on the initial design
and $3.265$\,GHz on the final one, $7.3\%$ and $8.7\%$ above the
published values; both resonator frequencies extrapolate to $2.7\%$ above
theirs
(Table~\ref{tab:crosscheck}). Every published value lies below its band.
With no extrapolation, the
highest uniform levels, $p = 2$ on $h/2$ ($7.0$ million unknowns on the
initial geometry; $p = 1$ on $h/4$ has more unknowns and lower values),
already lie $4.2\%$ and $4.7\%$ above the published qubit frequencies. The
estimates rest on Palace's $p = 1$ $h/4$ and $p = 2$ levels, which
geode-fem did not run. On the levels it did run, geode-fem and Palace
agree on the qubit to $10^{-4}\,\%$ on both geometries. On the resonator
they differ by $0.03\%$ at $h$, where Palace adds $4{,}946$ (initial) and
$5{,}944$ (final) elements by local bisection, and not at all in the six
archived decimals at $h/2$, where Palace's logs report no bisection and
the two codes differ only in the junction term
(Section~\ref{sec:junction}; the counts are in the logs). Three checks
test the estimates: a second refinement family and Palace's adaptive
refinement (Figure~\ref{fig:initialq}; the adaptive runs check only the
qubit), and a common-limit fit.

\paragraph{A second refinement family.}
Family~b refines the same $h$ mesh with Palace's own uniform refinement
instead of gmsh's split (family~a). It first bisects where a sheet is one
element wide (Section~\ref{sec:agreement}), so it is graded slightly
toward the sheets, a plausible reason for its higher observed rates,
$0.91$ (initial) and $0.90$ (final) against family~a's $0.62$ and $0.61$.
Its $h/2$ qubit is $3.9469$\,GHz against gmsh's $3.8771$\,GHz, and at
$h_F/2$ the two refinements of the fixture give $3.9528$ and $3.8939$\,GHz, so the
refinement scheme accounts for $1.5$--$1.8\%$ of the $h/2$ level and the
choice of starting mesh for only $0.15$--$0.43\%$. Its qubit estimates lie
$0.64\%$ and $0.83\%$ above family~a's, and every family-b band lies
inside the family-a band (Table~\ref{tab:crosscheck}). We therefore quote
the initial qubit frequency as $4.444$--$4.472$\,GHz, $7.3$--$8.0\%$ above the
published value, and the final one as $3.265$--$3.292$\,GHz, $+8.7\%$ to
$+9.6\%$ over the published $3.005$\,GHz.

\begin{figure}[t]
  \centering
  \anvilfig{fig3-initial-qubit.pdf}{fig3\_initial\_qubit.py}
  \caption{Checks of the extrapolated estimate on the initial geometry.
  (a, c) Every Palace level against the number of N\'ed\'elec unknowns, for
  the two refinement families of $h$ (filled: gmsh's split, as in
  Figure~\ref{fig:crosscheck}; open: Palace's) and Palace adaptive
  refinement from $h$ at $p = 2$ (diamonds, passes numbered in (a)); the
  dotted vertical line is the adaptive run's size limit, $10^7$ unknowns.
  (b, d) The estimates of the
  two families with their bands and, for the qubit only, the adaptive
  run's ill-conditioned fit (the resonator fit is unreliable). In (a, b), the
  circuit-model cap. Sources:
  \texttt{benchmarks/transmon\_devicelayout\_crosscheck/results.toml} and
  \texttt{benchmarks/transmon\_devicelayout\_controls/results.toml}.}
  \label{fig:initialq}
\end{figure}

\paragraph{Adaptive refinement.}
A further check comes from Palace's adaptive mesh refinement, run on each
geometry from its ladder's $h$ mesh at $p = 2$. Each pass solves, estimates
the error per element (summed over the qubit and resonator modes), and
refines the elements that carry the largest $70\%$ of it (D\"orfler
marking at $0.7$), without keeping the mesh conformal
(Figure~\ref{fig:initialq} for the initial geometry). On the initial
geometry, in five passes from $0.95$ to $11.0$ million unknowns, the qubit
frequency rises from $4.1683$ to $4.4297$\,GHz, $0.32\%$ below the estimate, inside
family a's band and below family b's; the third pass, at $2.2$ million unknowns, already exceeds the
highest uniform level ($4.3118$\,GHz at $7.0$ million). On the final
geometry, from $1.17$ to $14.6$ million unknowns, the qubit frequency rises from
$3.0164$ to $3.2468$\,GHz, $8.0\%$ above the published $3.005$\,GHz, and
the third pass again exceeds the highest uniform level ($3.1884$ against
$3.1464$\,GHz). The finest final pass lies $0.56\%$ and $1.39\%$ below the
two family estimates, again inside family a's band and below family b's. Both finest passes are computed
values, not extrapolations, and neither run reached a plateau: each stopped
when its last pass exceeded the run's size limit of $10^7$ unknowns, not on
its tolerance, and the last pass still moved the qubit by $+0.55\%$ (initial)
and $+0.63\%$ (final). A three-point fit $f_\infty - CN^{-\alpha}$ through
the last three passes gives $4.4552$\,GHz ($\alpha = 0.80$) and
$3.2672$\,GHz ($\alpha = 0.83$), but $+0.1\%$ on the finest pass moves
them by $24$ and $16$\,MHz, so they are ill-conditioned estimates. The
adaptive resonators rise more slowly and still lie below their bands, so
the adaptive runs check the qubit estimates only.

\paragraph{A common-limit fit.}
The last check fits the $p = 1$ ($h/2$, $h/4$) and $p = 2$ ($h$, $h/2$)
series of each ladder with one common limit and rate, exactly determined; it lies
$0.5$--$1.3\%$ above the rule's estimate on every family-a series, as
expected of a rule that leans low, and $+0.1\%$ on the $p = 2$ $h/2$ level
alone moves the initial-qubit fit by $20$\,MHz.

\section{A Static Route: A Bounded Island Capacitance}
\label{sec:static}

\paragraph{The bound.}
For the reduced electrostatic problem (island at potential 1, ground and
resonator at 0, feedline floating with zero net charge, and zero normal
flux on the vacuum box), the potential
$\phi$ satisfies $0 \le \phi \le 1$ by the maximum principle, and
$\Cisl = \int \varepsilon\,\nabla\phi\cdot\nabla\phi$ with $\varepsilon
= \varepsilon_0\varepsilon_r$. Every conforming finite-element energy is an
upper bound on $\Cisl$, because $\Cisl$ is the minimum of the energy
over admissible potentials. For the lower bound, take any piecewise
polynomial field $D$ on the mesh cut along the zero-thickness sheets.
Integrating by parts and using $0\le\phi\le1$ on every term gives
\begin{equation}
  \int D\cdot\nabla\phi \;\ge\; N(D) \;=\; Q_{\mathrm{island}}(D)
  - \lvert Q_{\mathrm{feed}}(D)\rvert
  - \lVert \operatorname{div} D\rVert_{L^1}
  - \lVert [D\cdot n]\rVert_{L^1(\mathcal F)}
  - \lVert D\cdot n\rVert_{L^1(\mathcal B)},
  \label{eq:cert-n}
\end{equation}
where $Q$ are the fluxes of $D$ into the two conductors (for the exact
field $D = \varepsilon\nabla\phi$, the negative of the physical
displacement, they are the conductor charges),
$[D\cdot n]$ is the normal jump on the interior faces $\mathcal F$ off
the sheets, and $D\cdot n$ on the box faces $\mathcal B$ is the one-sided
flux, which the zero-flux condition requires of the exact field but not
of $D$. Cauchy--Schwarz, $\bigl(\int D\cdot\nabla\phi\bigr)^2 \le
\int D\cdot\varepsilon^{-1}D \int\varepsilon\nabla\phi\cdot\nabla\phi$,
then gives, whenever $N(D) > 0$,
\begin{equation}
  \Cisl \;\ge\; \frac{N(D)^2}{\int D\cdot\varepsilon^{-1}D}.
  \label{eq:cert}
\end{equation}
The bound holds for any $D$, so solver tolerance can only weaken it. Each
$L^1$ norm is bounded above element by element and face by face,
$\lVert p\rVert_{L^1(K)} \le |K|^{1/2}\lVert p\rVert_{L^2(K)}$, with
$\lVert p\rVert_{L^2}^2$ integrated exactly by quadrature. We
take $D$ from the P2 potential, equilibrated patch by patch into
Raviart--Thomas fluxes of order one to three, with an independent
Crouzeix--Raviart flux as a check. The two bounds hold by theorem, for the
exact arithmetic of the discretization on the polygonal geometry of the
mesh, with lower bounds floored and upper bounds ceiled; the implementation
is checked by tests, not formally verified. The tests compare it with an
exact parallel plate and an analytic coaxial bracket, and mutation tests
reject a grounded-feedline or jump-free lower bound. The bound is evaluated
in plain f64, not interval arithmetic; the worst-case summation error is about
$10^{-10}$ relative, far below the rounding. Refinement does not move the
polygonal surfaces, so all meshes of one geometry bound the same exact
capacitance, and a family bracket is the largest lower bound and the
smallest upper bound.

\paragraph{Two box models.}
The bound above puts no condition on the box walls; this is the
\emph{natural} box, in which no field line leaves the box. The
eigenproblem instead makes the walls PEC (Section~\ref{sec:geometry}),
the \emph{grounded} box, with the walls held at ground potential.
Constraining the potential on the box,
to zero or to a floating constant, shrinks the set of admissible
potentials, so it can only raise $\Cisl$: the natural-box lower bound, and
the cap below that uses only it, hold for the grounded box too. For an upper
bound we re-solve the conforming P2 problems on the same meshes with the
box grounded.

\paragraph{The unmodified device.}
The \emph{unmodified} device is the initial design as published, before
the re-optimization of Section~\ref{sec:reopt} changes $c$, meshed as the
fixture. Uniform refinement of the fixture with P1 and P2 elements gives the upper
bounds of Figure~\ref{fig:convergence}a, from $137.71$ to $83.02$\,fF
(scalar) and from $136.69$ to $85.01$\,fF (tensor), ceiled. On adaptively
refined meshes the lower bound rises toward the falling upper bound
(Figure~\ref{fig:convergence}b): the island capacitance of the unmodified
initial design is bounded by $[77.40, 77.61]$\,fF with the scalar
$\bar\varepsilon$ and by $[79.56, 79.76]$\,fF with the sapphire tensor.
The bounded tensor-minus-scalar
difference, computed from the unrounded bounds, is $1.95$--$2.35$\,fF, or
$2.5$--$3.1\%$ of the scalar value. With the box grounded, the tensor
$\Cisl$ is bounded by $[79.56, 79.77]$\,fF. On the finest mesh, grounding
the box raises the conforming value by $0.0065$\,fF, $8\times10^{-5}$
relative, an estimate of the effect of the box model, whose size is
bounded only by $[0, 0.21]$\,fF. The admissible sets are nested, so the
grounded bracket also covers a floating box. With $C_J$ added, the transmon's
$C_\Sigma = \Cisl + C_J$ is bounded by $[85.06, 85.27]$\,fF for the
eigenproblem's grounded box (tensor; $[85.06, 85.26]$\,fF natural). The
fixture itself overstates the tensor $\Cisl$ by at least $56.92$\,fF, which is at least $71.3\%$ of
the exact value, and the scalar one by at least $77.4\%$.

\begin{figure}[htbp]
  \centering
  \anvilfig{fig4-convergence.pdf}{fig4\_convergence.py}
  \caption{Mesh convergence of the island capacitance $\Cisl$ of the
  unmodified device, with the scalar trace-averaged
  permittivity (blue) and the sapphire tensor (orange). (a) Uniform
  refinement of the fixture ($h_F$, $h_F/2$, $h_F/4$): every marker is a
  conforming upper bound on the exact capacitance of the meshed polygonal
  geometry. The short bars are the two-sided bounds from
  adapted meshes, each drawn at the tetrahedron count of the mesh that
  sets its lower bound. (b) The adaptive passes behind those bars: on each
  adapted mesh the conforming P2 value (hollow, an upper bound) and the
  best equilibrated-flux lower bound (filled) close in on each other; the
  title gives the bounded tensor-minus-scalar gap. All values are for the
  natural box; the grounded box changes only the tensor upper bound, to
  $79.77$\,fF. Sources: \texttt{c\_sigma\_convergence.toml},
  \texttt{c\_sigma\_bracket.toml} and
  \texttt{c\_sigma\_bracket\_tensor.toml}.}
  \label{fig:convergence}
\end{figure}

\paragraph{From capacitance to the qubit frequency.}
Reading the grounded-box tensor bound through the two-node circuit of
Section~\ref{sec:junction}, over every corner of the $\Cisl$ bound and the
$L_s$ and $k$ ranges, with the resonator at its extrapolated estimate, gives the
initial dressed qubit in $[4.449, 4.458]$\,GHz; the natural-box bound raises
its low end by about $0.3$\,MHz. That range lies inside the eigenmode band of
Section~\ref{sec:crosscheck} and inside the range of the two family
estimates, and family~a's estimate is $0.22\%$ below its midpoint, a
consistency check (Section~\ref{sec:intro}). The range is an estimate, because $L_s$, $k$
and the resonator estimate are.

\paragraph{A model-dependent cap.}
Under the circuit model, a series inductance and a coupling to a higher
mode can only lower the qubit, so the lower bound $\Cisl \ge 79.56$\,fF,
that is $C_\Sigma \ge 85.06$\,fF, caps it at $1/(2\pi\sqrt{L_J C_\Sigma})
= 4.477$\,GHz. The cap is rigorous only given that model (one island node
with $L_J \parallel C_J$, $L_s \ge 0$, coupling only to modes above the
qubit, a quasi-static $C_\Sigma$) and rests on the fixture, not on the
ladder meshes, which agree with it to $0.02\%$ at $h$, so it is
model-dependent and not in the Bounded tier of Table~\ref{tab:tiers}. It
does not depend on the
extrapolation rule. Every computed level and both family estimates lie
below it; family~a's common-limit fit, $4.4802$\,GHz, exceeds it by
$3.6$\,MHz, within that fit's conditioning. With no extrapolation, the finest adaptive pass and the cap put the initial qubit between $4.4297$ and
$4.477$\,GHz, $+7.0\%$ to $+8.1\%$ over the published value; the lower end is a computed value, the upper end holds only under the circuit model.

\section{Re-Optimization and Its Controls}
\label{sec:reopt}

At the extrapolated frequencies the published final design misses its
targets. We repeat DeviceLayout's optimization against its own
discretization, to compare optimizers in one environment, and against the
extrapolated frequencies, to find the design that meets the targets.
Cost is counted in \emph{Palace runs}: one Palace eigensolve at
DeviceLayout's discretization, about $4$\,min on 32 ranks. A Newton
\emph{step} places a new design and costs one Palace run; a \emph{probe}
is a Palace run that is not a step. geode-fem solves are not Palace runs,
and their cost is given in seconds. Newton's Jacobian has one column per
design parameter, $c$ and $L$; the next two paragraphs give how each is
obtained.

\paragraph{The capacitor\_length derivative.}
DeviceLayout builds the island upward from the fixed junction, so a change
of $c$ stretches the strip that holds the island below its top and
translates every node above that strip (Figure~\ref{fig:geometry}a).
geode-fem turns this motion into a node velocity
$\partial X/\partial c$, checks that every metal outline node in the
stretched strip moves along its own edge ($3{,}146$ nodes on the initial
mesh), and evaluates Eq.~\eqref{eq:hf-eig} for both modes in $0.44$\,s.
On the initial mesh it matches a central
difference of two re-solves of the moved mesh ($\pm0.5\,\mu$m) to every
printed digit: $-2.443227\times10^{-3}$\,GHz/$\mu$m for the qubit and
$-3.964\times10^{-6}$\,GHz/$\mu$m for the resonator.

\paragraph{The total\_length column.}
DeviceLayout re-lays the meander when $L$ changes, so no fixed-topology
node motion represents it. A transmission-line model holds $f_r(L+\delta)$
constant and takes $\partial f_q/\partial L = 0$, with $\delta =
216\,\mu$m from one finite difference of two geode-fem $p = 1$ solves
($L = 5000$ and $5250\,\mu$m). $\delta$ is weakly determined: remesh noise
of about $0.15\%$ in either $f_r$ moves it by about $\pm175\,\mu$m. The
model enters only the Jacobian, so a poorer Jacobian costs steps, not
accuracy.

\paragraph{Newton on DeviceLayout's objective.}
DeviceLayout re-meshes every design, so the $p = 1$ objective on $h$ is not
smooth in the parameters: from $c = 620$ to $650\,\mu$m the remeshed
$p = 1$ difference of the qubit is $-5.23\times10^{-3}$\,GHz/$\mu$m
against a fixed-topology derivative of $-2.28\times10^{-3}$ averaged over
the interval (the mean of its end values), because remeshing alone moves
the qubit by about $2.6\%$, more than the $2.0\%$ the step moves it. The
iteration therefore runs Newton directly on the objective $F$, the
frequencies at DeviceLayout's discretization,
in DeviceLayout's variables $(x_1, x_2)$, in which the frequencies are
close to linear and nearly decoupled: a probe in either variable moves its
own mode by $4.5\%$ and the other by at most $0.05\%$. The $c$ column of
the Jacobian is geode-fem's $p = 1$ derivative on $h$, $\partial f_h/\partial c$, scaled by the
ratio $S = F/f_h$ of the two discretizations at the same design. Each step costs one geode-fem
$p = 1$ solve with its gradient ($19$--$25$\,s) and one Palace run; the
geode-fem solves of the one-time $\delta$ calibration are not Palace runs.

\begin{table}[htbp]
  \centering
  \caption{Palace runs to DeviceLayout's criterion ($\mathrm{err} <
  10^{-4}$) at its discretization, all in our Palace environment except the
  published PRIMA run. Every Palace run a method uses counts, with the
  initial design as run~1 and probes included, among them the probe that
  sets the Palace $\delta$ (Broyden's second probe, shared by the two
  rows); geode-fem solves are not counted. Each count is one trajectory.
  PRIMA runs at \texttt{rhobeg} $= 0.2$, its initial trust-region radius,
  unless stated; the Newton variants
  use the transmission-line $L$ column.
  Sources: the \texttt{optimize} and \texttt{controls} artifacts under
  \texttt{benchmarks/transmon\_devicelayout\_*/}.}
  \label{tab:controls}
  \footnotesize
  \setlength{\tabcolsep}{4pt}
  \begin{tabular}{@{}lllccc@{}}
    \toprule
    Method & $c$ column & $\delta$ from & Palace runs & err after first step & Design at criterion ($\mu$m) \\
    \midrule
    PRIMA, published & none & none & 9 & unpublished & (1217.678, 7243.955) \\
    PRIMA, our environment & none & none & 10 & n/a & (1209.125, 7207.551) \\
    \quad \texttt{rhobeg} $0.1$ & none & none & 10 & n/a & (1235.705, 7216.766) \\
    \quad \texttt{rhobeg} $(x_u - x_l)/2$ & none & none & 12 & n/a & (1245.900, 7291.401) \\
    Newton & shape derivative & geode-fem & 3 & $7.9\times10^{-4}$ & (1225.346, 7223.250) \\
    Newton & scaling law & geode-fem & 3 & $1.7\times10^{-4}$ & (1232.664, 7221.270) \\
    Newton & scaling law & Palace & 4 & $2.1\times10^{-4}$ & (1229.026, 7225.712) \\
    Broyden & two probes & none & 5 & $6.2\times10^{-4}$ & (1224.847, 7232.447) \\
    \bottomrule
  \end{tabular}
\end{table}

\begin{figure}[!tbp]
  \centering
  \anvilfig{fig5-reoptimization.pdf}{fig5\_reoptimization.py}
  \caption{The re-optimization and its controls. (a) The design plane near
  the optimum, in two regions. D: the published final design, PRIMA's end
  point in our environment and the Newton iterates against DeviceLayout's
  discretization (the first legs start at the initial design, off the
  plot). E: the steps against the extrapolated frequencies, with the
  step-2 qubit band
  mapped linearly to capacitor\_length (an estimate); the text gives the
  published and re-optimized capacitor lengths. (b) DeviceLayout's
  err against Palace runs at its discretization in our environment,
  probes counted (the Palace-$\delta$ Newton's run 2 is the $L$ probe it
  shares with Broyden), PRIMA at \texttt{rhobeg} $= 0.2$; the published PRIMA
  point is at its 9 runs and the err of its published frequencies,
  $1.151\times10^{-5}$ (DeviceLayout's environment). This counts Palace
  runs, not total cost: each shape-derivative Newton point also
  costs one geode-fem solve. (c, d) Extrapolated
  estimates and bands against the targets, with step~2 in both refinement
  families. Sources: the \texttt{optimize}, \texttt{controls} and
  \texttt{crosscheck} artifacts under
  \texttt{benchmarks/transmon\_devicelayout\_*/}.}
  \label{fig:reopt}
\end{figure}

\paragraph{Against DeviceLayout's discretization.}
Table~\ref{tab:controls} and Figure~\ref{fig:reopt}(b) compare the
optimizers on one objective. Newton
with the shape-derivative column meets the criterion at its second step,
after 3 Palace runs, and its three geode-fem solves add about $70$\,s,
against about $4$\,min per Palace run; a third step reaches
$\mathrm{err} = 1.4\times10^{-7}$ at $(1236.360, 7224.900)\,\mu$m. Two
gradient-free controls test where the 3 runs come from. The first replaces
the $c$ column by the scaling law $\partial F_q/\partial x_1 = F_q/x_1$,
which follows from $f_q \propto 1/\sqrt{C_\Sigma}$ with $C_\Sigma$ taken
proportional to $c$. With the
geode-fem $\delta$ it also takes 3 runs, and its first step has an err
$4.7$ times lower (Table~\ref{tab:controls}), about $2.2$ times closer in
frequency. With $\delta = 563.879\,\mu$m taken
instead from a Palace probe ($c$ fixed, $L$ from $5000$ to
$5263\,\mu$m), so that no geode-fem input remains, it takes 4 runs: the
same two steps, plus the probe. This $\delta$ is $2.6$ times the geode-fem
one, more than the $\pm175\,\mu$m noise estimate, and costs no extra step.
The second control, Broyden's method \citep{broyden1965class} from two
forward finite-difference probes, takes 5.

\paragraph{PRIMA in our environment.}
DeviceLayout's own PRIMA call,
re-run verbatim in our environment, takes 10: runs 1--5 are its initial
interpolation set of $2n+1 = 5$ points \citep{powell2009bobyqa}, and run 9
misses the criterion by $15\%$ ($\mathrm{err} = 1.150\times10^{-4}$). In
this environment the published final design already meets the criterion
($\mathrm{err} = 1.5\times10^{-5}$), so the environment moves PRIMA's
trajectory, not the attainability of the target; the published count of 9
is from DeviceLayout's environment. With \texttt{rhobeg}
at $0.1$ instead of $0.2$, PRIMA also takes 10 runs, and
at the largest radius BOBYQA accepts on DeviceLayout's box, $(x_u -
x_l)/2$, it takes 12. Each setting is a single run. In one environment, then,
Newton's method in DeviceLayout's variables takes 3 or 4 Palace runs with
any of the three Jacobians of Table~\ref{tab:controls} (the shape
derivative, and the scaling law with either $\delta$) and PRIMA 10 to 12;
the one run between 3 and 4 is the Palace probe that geode-fem's $\delta$
calibration replaces.

\paragraph{Where the first steps land.}
The shape-derivative and Broyden first steps land at $c = 1338.338$ and
$1320.522\,\mu$m, past about $1230\,\mu$m at the criterion: both follow
a local slope, $\partial F_q/\partial x_1 = 3.66$ and $3.71$\,GHz, within about $1.5\%$ of each other and $11$--$12\%$ shallower than the scaling law's $F_q/x_1 = 4.168$\,GHz.

\paragraph{Against the extrapolated frequencies.}
Starting from the published final design, each Newton step uses the
extrapolated estimate and the geode-fem derivative at the current design,
and each new design is evaluated by a full ladder with the rule of
Section~\ref{sec:crosscheck} (Table~\ref{tab:newton-conv}).
The first overshoots the qubit by $-1.09\%$; the second lands at $+0.16\%$
and $-0.04\%$ of the targets, with both targets inside the bands. On the
ladders of both steps, geode-fem and Palace agree on the qubit to
$2\times10^{-4}\,\%$.
One scaling-law step from the same start, evaluated by the Palace half of
the ladder, lands $3.6$ and $8.2\,\mu$m short of step~2 and at $+0.13\%$
and $+0.12\%$ of the targets: one gradient-free step does what the two
shape-derivative steps did.

\begin{table}[htbp]
  \centering
  \caption{Steps against the extrapolated frequencies, in GHz. Each design
  is evaluated by a full resolution ladder (the scaling-law step and
  family b by Palace alone) and the rule of
  Section~\ref{sec:crosscheck}; values are estimates with their bands.
  Family b refines step 2's $h$ mesh with Palace's uniform refinement
  instead of gmsh's. Sources: the \texttt{optimize}, \texttt{controls} and
  (published final) \texttt{crosscheck} artifacts under
  \texttt{benchmarks/transmon\_devicelayout\_*/}.}
  \label{tab:newton-conv}
  \footnotesize
  \begin{tabular}{@{}lcccc@{}}
    \toprule
    Design & $c$ ($\mu$m) & $L$ ($\mu$m) & $f_q$ [band] & $f_r$ [band] \\
    \midrule
    Published final & 1217.678 & 7243.955 & 3.265 $[3.236, 3.471]$ & 4.088 $[4.077, 4.139]$ \\
    Step 1 & 1481.657 & 7407.945 & 2.967 $[2.938, 3.188]$ & 4.004 $[3.993, 4.055]$ \\
    Step 2 & 1446.078 & 7416.352 & 3.005 $[2.979, 3.221]$ & 3.998 $[3.987, 4.067]$ \\
    Step 2, family b & 1446.078 & 7416.352 & 3.020 $[3.003, 3.101]$ & 4.022 $[4.016, 4.043]$ \\
    Scaling-law step & 1442.479 & 7408.181 & 3.004 $[2.977, 3.235]$ & 4.005 $[3.994, 4.055]$ \\
    \bottomrule
  \end{tabular}
\end{table}

\paragraph{The re-optimized design.}
Step 2 is the design whose extrapolated frequencies sit on the targets:
$c = 1446.078\,\mu$m, $18.8\%$ longer than published, with a resonator
$2.4\%$ longer ($+172\,\mu$m). Family~b at step 2 gives
$3.0203$ $[3.0033, 3.1014]$\,GHz for the qubit and $4.0225$\,GHz for the
resonator, $0.52\%$ and $0.60\%$ above family~a, the same spread as
on the initial geometry; both targets lie just below its band, so on family
b step 2 is slightly short. The headline length we quote is not this
design's $c$ itself but the two families' qubit estimates mapped to
$3.0$\,GHz by the scaling law $f_q \propto c^{-1/2}$: they call for $c =
1451$ and $1466\,\mu$m, $+19.1\%$ and $+20.4\%$ over the published length. Their
bands map to $+17\%$ to $+37\%$ (family a) and $+19\%$ to $+27\%$
(family b) under the scaling law, and to $+17\%$ to $+39\%$ ($c$ from about
$1423$ to $1690\,\mu$m) and $+19\%$ to $+28\%$ linearly through $dF/dc
\approx S\,\partial f_h/\partial c = -9.05\times10^{-4}$\,GHz/$\mu$m at
the step-1 design, centred on the target. Under either mapping the first
range contains the second. Since the estimates lean low, the common-limit
fit, which puts step~2 at $3.039$ and $4.019$\,GHz, would call for a longer
capacitor still.

\section{Discussion}
\label{sec:discussion}

\paragraph{What is established, and how.}
Table~\ref{tab:tiers} sorts the results by the kind of evidence behind
them. The brackets are the only bounds, and the cap follows from a bracket
only under a circuit model.

\begin{table}[htbp]
  \centering
  \caption{Evidence tiers. Bounded: a two-sided bound that holds by theorem
  for the meshed polygonal geometry. Model-dependent: a rigorous consequence of a
  bound under a stated circuit model. Computed: a solve on a stated
  discretization. Estimated: an extrapolation with a stated rule.
  Consistency: an agreement that is not independent evidence.}
  \label{tab:tiers}
  \footnotesize
  \begin{tabular}{@{}>{\raggedright\arraybackslash}p{0.17\linewidth}>{\raggedright\arraybackslash}p{0.77\linewidth}@{}}
    \toprule
    Tier & Results \\
    \midrule
    Bounded & Island capacitance $\Cisl$ of the unmodified device, without $C_J$, natural box: $[77.40, 77.61]$\,fF (scalar), $[79.56, 79.76]$\,fF (tensor); grounded box, the eigenproblem's model: $[79.56, 79.77]$\,fF (tensor), so $C_\Sigma = \Cisl + C_J$ in $[85.06, 85.27]$\,fF and $e^2/2hC_\Sigma$ in $[0.2271, 0.2278]$\,GHz; tensor-minus-scalar gap $2.5$--$3.1\%$ \\
    Model-dependent & The initial-qubit cap of $4.477$\,GHz $= 1/(2\pi\sqrt{L_J C_\Sigma})$ at the lower bound on $C_\Sigma$: two-node circuit, quasi-static $C_\Sigma$, bounded on the fixture's polygonal geometry, not on the ladder meshes; with the finest adaptive pass, $[4.4297, 4.477]$\,GHz; the same $e^2/2hC_\Sigma$ read as the transmon's charging energy $E_C/h$ \\
    Computed & Same-mesh agreement (Table~\ref{tab:agreement}); the reproduction within $0.7\%$; every ladder level and adaptive pass (adaptive qubit frequencies $+7.0\%$ and $+8.0\%$, still rising); the gradient check; Palace runs to DeviceLayout's criterion, one trajectory each (Table~\ref{tab:controls}) \\
    Estimated & Extrapolated frequencies and bands (Table~\ref{tab:crosscheck}); the adaptive and common-limit fits; the static-route qubit range; $L_s$ and $k$; the box-model effect on $\Cisl$, $8\times10^{-5}$; the re-optimized capacitor, $+19$ to $+20\%$ over two families and about $+17$ to $+39\%$ within the band \\
    Consistency & The two routes agree to $0.22\%$ on the initial qubit; the extrapolation rule was fixed after the static route was known, so this is not a blind prediction \\
    Not claimed & Agreement with measurement; a verified hit of $3.0$/$4.0$\,GHz in the $h \to 0$ limit; a global optimum; an advantage of the shape derivative on this problem; eigenmode-participation (EPR) derivatives, which need eigenvector derivatives \citep{nelson1976simplified}; a speed advantage over Palace \\
    \bottomrule
  \end{tabular}
\end{table}

\paragraph{Where the shift comes from.}
The evidence places the $7$--$8\%$ shift in the discretization, not in
arithmetic precision or in the eigensolver. geode-fem computes in double
precision (f64), and Palace is the build pinned in
\path{reference/palace/docker/Dockerfile}, which sets no single-precision
option. On the levels both codes ran, the two coarsest family-a levels
with linear elements, the two codes agree on the qubit frequency to
$2\times10^{-4}\%$, the looser of the ladder and re-optimization values
(Sections~\ref{sec:crosscheck} and~\ref{sec:reopt}); on those levels, a
loss of precision large enough to matter in either code would appear as a
same-mesh disagreement far above that level. Palace is used here as a
frequency-domain finite-element eigensolver, so there is no time stepping
or numerical dispersion to suspect: it returns the eigenmodes of the
discrete problem that its input
mesh and element order define, and given the same mesh, geode-fem, with
its own assembly and eigensolver, returns the same modes
(Table~\ref{tab:agreement}; the crosses in Figure~\ref{fig:crosscheck}).
The error lies in that input: refining the mesh from $h$ to $h/2$ moves
both codes' frequencies up together (Figure~\ref{fig:crosscheck}), and
the published values come from one discretization.

\paragraph{Why the qubit is approached from below.}
The qubit mode is quasi-static; the quasi-static part of a first-order
N\'ed\'elec field is the gradient of a P1 potential, which can only
overstate the capacitance, so the computed qubit should approach its limit
from below: on the fixture the P1 bound ($136.69$\,fF, with $C_J$ added)
overstates $C_\Sigma$ by a factor of $1.67$, and the family estimates exceed
the fixture's $p = 1$ qubit by $1.66$--$1.68$ in $f^2$. This is a heuristic, not a proof: it omits $L_s$,
the coupling, the rotational field and the distributed resonator.

\paragraph{Why the shape derivative gives no advantage here.}
Three properties of this benchmark explain why Newton's method needs so
few Palace runs whatever its Jacobian. First, the problem has
two nearly decoupled parameters, close to linear in DeviceLayout's
variables, so Newton's method converges with any reasonable Jacobian: the
shape derivative or the scaling law $f_q \propto x_1$, which costs nothing,
in 3 or 4 runs (Table~\ref{tab:controls}), and even Broyden's two-probe
secant in one or two runs more. Second, a finite-difference gradient costs only $N+1 =
3$ solves. Third, the coarse-mesh, fixed-topology derivative is the slope
of a different discrete problem from the one the iteration evaluates
(DeviceLayout remeshes every design and Palace evaluates it on a curved
mesh at $p = 2$); once scaled by the ratio $S$ it gives about what
Broyden's probe gives (Section~\ref{sec:reopt}). That the derivative is exact for its own
discrete problem and cheap ($0.44$\,s) does not change the number of
Palace runs.

Derivatives pay off in other regimes. With $N$ parameters, a discrete
adjoint gives all $N$ derivatives of one objective for about one extra
solve, where finite differences need $N$ more; coupled targets with no
per-parameter scaling law need a computed Jacobian; and freeform shapes
have many parameters by construction. A companion paper on adjoint shape
optimization of a conformal patch antenna \citep{walters2026conformal}
measures this on the worst-of-band reflection $|S_{11}|$ over
$2.400$--$2.4835$\,GHz (artifact:
\texttt{benchmarks/patch\_antenna\_conformal/tuned/tuned\_summary.toml}).
At 38 shape parameters the adjoint reaches a $-9.5$\,dB target with about
$11$ times fewer solves than the fastest gradient-free run that reached it,
a tuned full-covariance CMA-ES comes within $7.7$ times at $-9$\,dB, and
runs that exhausted their budget before reaching the $-9.5$\,dB target
bound the ratio from below at $2.2$--$3$; at 4 parameters the ratio is $1.1$ at
$-9$\,dB and $2.4$ at $-9.5$\,dB.

\paragraph{Threats and limitations.}
Agreement between two codes is evidence against implementation error, not
against shared modelling choices such as the port sheet behind the
$17.49$\,GHz mode. The extrapolations assume the rate the sheet-edge
singularity allows on uniform meshes, their high sides are wide, neither
adaptive run reached a plateau, and $L_s$ and $k$ were extrapolated from
two or three discretizations. Every limit $h \to 0$ here is a limit of
DeviceLayout's model (zero-thickness PEC sheets, lossless and linear, open
readout ports in the ladders); film thickness and kinetic inductance remove
the edge singularity and move the frequencies by amounts not measured
here. The element edges
along the sheet outlines are mostly $10$--$12\,\mu$m long on $h$ and so
$2.5$--$3\,\mu$m on $h/4$, more than ten times a $0.2\,\mu$m film, on
every uniform mesh. Palace's adaptive runs refine at most four times: a
typical outline edge stays above about $0.6\,\mu$m, but the
$1.0$--$1.8\,\mu$m edges of the junction leads, where the qubit's error
indicator is largest, could reach $0.06$--$0.11\,\mu$m, comparable to or
below the film (an estimate: no committed record holds the adaptive
meshes' minimum edge). There the zero-thickness model, which every mesh
discretizes, is an approximation whose effect is not quantified; the shift
measured here is the discretization error within that model, for one
device.

\paragraph{What the shift means for a design.}
Table~\ref{tab:design} reads the initial qubit as a total capacitance. The
bound gives $E_C/h = e^2/2hC_\Sigma$ in $[0.2271, 0.2278]$\,GHz for the
grounded box, Bounded as a function of $C_\Sigma$ and model-dependent as
the transmon's charging energy (a single-island lumped picture). Under the
two-node circuit with the $L_s$ and $k$ ranges, the published $4.14$\,GHz
implies a $C_\Sigma$ $15.5$--$16.0\%$ above the bound, so $E_C$ and the
anharmonicity, close to $-E_C$ \citep{koch2007charge}, are larger than the
published discretization implies by as much. The bare LC reading at fixed
$L_J$ omits $L_s$ and $k$, so it reads every frequency as a larger
$C_\Sigma$, $16.6$--$17.0\%$ above the bound for the published value
(rounded outward). Of that excess, taken as a factor, the computed finest adaptive pass accounts
for $1.145$, and the remaining $1.9$--$2.1\%$ splits between the still-rising pass and $L_s$ and $k$. The transition $f_{01}$ rises by a
little less than the linear mode, about $6.6\%$ at the finest adaptive pass, and
$L_J$ can retune the qubit frequency but not $E_C$.

\begin{table}[htbp]
  \centering
  \caption{The initial qubit read as a total capacitance $C_\Sigma$, with
  $E_C/h = e^2/2hC_\Sigma$; ``bare'' is the LC reading at fixed $L_J$
  ($L_s = k = 0$), ``two-node'' the circuit of
  Section~\ref{sec:junction} with the $L_s$ and $k$ ranges, given first for
  the published value; bare rows omit $L_s$ and $k$, so their $C_\Sigma$
  lies above the (grounded-box) bound. Sources:
  \texttt{c\_sigma\_bracket\_grounded\_box.toml}, Table~\ref{tab:crosscheck}
  and \texttt{two\_node\_reading.toml} in the paper's directory (the
  \texttt{transmon\_loop\_inductance} anchor uses another topology).}
  \label{tab:design}
  \footnotesize
  \begin{tabular}{@{}llccc@{}}
    \toprule
    Source & Tier & $f_q$ (GHz) & $C_\Sigma$ (fF) & $E_C/h$ (GHz) \\
    \midrule
    Bound & Bounded & -- & $[85.06, 85.27]$ & $[0.2271, 0.2278]$ \\
    Static route & Estimated & $[4.449, 4.458]$ & $85.8$--$86.1$ (bare) & $0.225$--$0.226$ (bare) \\
    Adaptive finest pass & Computed & $4.4297$ & $86.87$ (bare) & $0.223$ (bare) \\
    Published & -- & $4.14$ & $[98.51, 98.64]$; $99.45$ (bare) & $0.196$--$0.197$; $0.195$ (bare) \\
    \bottomrule
  \end{tabular}
\end{table}

\paragraph{What a practitioner can do.}
The standard remedy for the slow edge rate is to refine toward the edges,
with a mesh graded toward the edges or with $hp$ refinement
\citep{demkowicz2006computing}: Palace's adaptive runs
(Section~\ref{sec:crosscheck}) and the gap-refined, fourth-order setting of
\citet{ye2025electromagnetic}, within $0.3\%$ of their extrapolated
values, are instances.
A first check for a copied layout is the strip width: on the fixture, $993$ of the $1{,}007$
edges that trigger Palace's bisection lie where a resonator or feedline
strip is one element wide (Section~\ref{sec:agreement}). Surface-loss estimates \citep{wang2015surface,martinis2022surface} need
the same edge resolution as the frequencies.

\section{Conclusion}
\label{sec:conclusion}

Palace reproduces the published DeviceLayout.jl SingleTransmon
frequencies at the example's own discretization, and that discretization
reads the qubit frequencies low: adaptive runs put them $7.0\%$ and $8.0\%$
above their published values with no extrapolation, and the estimates of
two refinement families, which lean low, put them $7.3$--$9.6\%$ above. geode-fem
reproduces Palace's qubit on every level both codes ran, and a two-sided
bracket on the island capacitance, read through a circuit model, is
consistent with that shift and caps the initial qubit under that model.
At the extrapolated frequencies the final design, to reach its $3.0$ and
$4.0$\,GHz targets, needs a capacitor about $19$--$20\%$ longer than
published. Against the example's own objective, Newton's method in DeviceLayout's
variables reaches the criterion in 3 or 4 Palace runs with any of three
Jacobians, against 10 to 12 for PRIMA; geode-fem's shape derivative gives
no advantage here. For design examples that are meant to be
copied, the discretization is part of the result and belongs in the report
beside it.

\section{Reproducibility}
\label{sec:repro}

Every result is read from an archived, scripted artifact.\footnote{The geode-fem source is public at
\url{https://github.com/rjwalters/geode-fem}. This paper's PDF and \LaTeX{}
source are archived at
\href{https://doi.org/10.5281/zenodo.23580545}{doi:10.5281/zenodo.23580545}.}
Each caption names its artifact; the Palace cells are under
\path{reference/fixtures/transmon_palace/}, and the DeviceLayout meshes are
generated by \path{reference/julia/generate_transmon_param.jl}. The versions are geode-fem commit
\texttt{3174015} for the same-mesh agreement, with the commit of each later
result recorded in its artifact, Palace \texttt{fba6a5b}, DeviceLayout.jl
v1.15.0, PRIMA.jl 0.2.4 and Burn 0.21. Palace ran on AWS
Graviton4 hosts, and the capacitance brackets on a 28-core workstation and
on Graviton4. geode-fem's
direct solves at $h/2$ peaked at $94$ and $136$\,GB, so only Palace ran
larger cells.

Palace's ladder cells use SLEPc shift-invert at tolerance $10^{-8}$,
\texttt{Target} $3.0$ (initial) or $2.0$\,GHz (final), and its default
direct solver, SuperLU ($p = 1$ on $h$), or AMS-preconditioned GMRES at relative
tolerance $10^{-9}$, on 16 ($p = 1$ on $h$) or 32 MPI ranks; the fixture ran
on 8, and the runs at DeviceLayout's discretization keep the example's
tolerances ($10^{-6}$ and $10^{-7}$, \texttt{Target} $2.5$\,GHz). geode-fem's Lanczos shifts are
$5$ and $4$\,GHz, with relative residuals below $4\times10^{-8}$.
The $h$, $h/2$ and $h/4$ meshes have $133{,}355$, $1{,}066{,}840$ and
$8{,}534{,}720$ tetrahedra (initial) and $163{,}431$, $1{,}307{,}448$ and
$10{,}459{,}584$ (final). The meshes were made
with gmsh 4.15.2 (\texttt{gmsh\_jll}) under Julia 1.10.11, and Palace was
built against the MPICH of the Rocky Linux 9 image in
\path{reference/palace/docker/Dockerfile}.

\paragraph{Acknowledgments.}
We thank the Palace and DeviceLayout.jl developers for publishing the
SingleTransmon example and its frequencies. Section~\ref{sec:built} gives
the author contributions and the AI-orchestrated development.

\section{Author Contributions and How geode-fem Was Built}
\label{sec:built}

geode-fem realizes the whiteroom operator specification
\citep{dunnavant2026whiteroom}, whose four layers, L1--L4,
restate Palace's cited source as pure functions, unfold those into
compositions of base algebraic primitives, recast them where possible as
global tensor operations, and express them in a small graph-evaluation
calculus. Author contributions: C.D. wrote that specification, gave design input and review on the solver's operator layer, and contributed code: the
migration to Rust 2024, the Metal (Apple GPU) backend, the present layout
of the core crate and the examples, and the geode-util crate. R.W. ran
the research and the AI agents that produced the rest of the code, the
experiments and the text of this paper, made every scoping and claim
decision, and is responsible for the results. The agents were AI models
working in separate roles: Anthropic's Claude (Opus, Fable and Sonnet
models, through Claude Code) and OpenAI's GPT Sol and GPT Astra models
(via Codex). They were coordinated by Loom \citep{loom}, an
issue, worktree and pull-request workflow, and by Squad, a message room
the agents share.
Nearly every change was reviewed by a separate agent session. The agent
reviews of this paper recomputed its numbers from the archived artifacts,
reproduced runs, and mutation-tested the scripts that check the text and
figures against them.

\bibliographystyle{plainnat}
\bibliography{refs}

\end{document}
